% anhang_e_glossar.tex
\chapter{Glossar}
\label{chap:anhang_e}

\section*{Einleitung}
Dieses Glossar erläutert die zentralen Begriffe, Konzepte und Abkürzungen, die in der GWI-D-Analyse verwendet werden. Es dient als Referenz für das Verständnis der methodischen Grundlagen, analytischen Konzepte und politischen Schlussfolgerungen.

\section{Grundbegriffe und Konzepte}

\begin{description}
    \item[Gemeinwohl-Index Deutschland (GWI-D)] \hfill \\
    Ein multidimensionaler Wohlfahrtsindex, der das Gemeinwohl in Deutschland auf Basis von vier Komponenten (Wirtschaft, Gerechtigkeit, Zukunft, Zusammenhalt) für den Zeitraum 1960–2023 misst.
    
    \item[Gemeinwohl] \hfill \\
    Ein normatives Konzept, das über individuellen Nutzen hinausgeht und das Wohlergehen der gesamten Gesellschaft sowie zukünftiger Generationen umfasst. Im GWI-D operationalisiert durch vier gleichgewichtete Dimensionen.
    
    \item[Multidimensionalität] \hfill \\
    Die Anerkennung, dass Wohlfahrt nicht durch eine einzige Dimension (wie BIP) erfasst werden kann, sondern mehrere, teilweise unabhängige Dimensionen umfasst.
    
    \item[Polykrise] \hfill \\
    Die gleichzeitige Überlagerung mehrerer, sich gegenseitig verstärkender Krisen (z.B. Gesundheits-, Klima-, Wirtschafts-, Demokratiekrise), wie sie insbesondere ab 2020 beobachtet wurde.
    
    \item[Stagflation] \hfill \\
    Eine wirtschaftliche Situation, in der Stagnation (kein oder negatives Wachstum) mit hoher Inflation einhergeht. Charakteristisch für die Periode 1974–1982 in Deutschland.
    
    \item[Agenda 2010] \hfill \\
    Ein umfassendes Reformpaket der rot-grünen Bundesregierung (2003–2005) mit tiefgreifenden Änderungen in Arbeitsmarkt-, Sozial- und Steuerpolitik.
    
    \item[Prekarisierung] \hfill \\
    Der Prozess der Zunahme unsicherer, schlecht bezahlter und sozial wenig abgesicherter Beschäftigungsverhältnisse, insbesondere seit den 1990er Jahren.
    
    \item[Resilienz] \hfill \\
    Die Fähigkeit von Individuen, Gesellschaften oder Systemen, Krisen und Schocks zu bewältigen und sich anzupassen, ohne dauerhafte Schäden zu erleiden.
    
    \item[Generationengerechtigkeit] \hfill \\
    Das Prinzip, dass politische Entscheidungen die Interessen und Lebenschancen zukünftiger Generationen berücksichtigen müssen, insbesondere in Bezug auf Umwelt, Staatsverschuldung und Sozialsysteme.
    
    \item[Sozialer Zusammenhalt] \hfill \\
    Der Grad der Verbundenheit und Solidarität in einer Gesellschaft, gemessen durch Vertrauen, Teilhabe, Integration und gemeinsame Werte.
\end{description}

\section{GWI-D-spezifische Begriffe}

\begin{description}
    \item[GWI-D-Komponenten] \hfill \\
    Die vier Hauptdimensionen des Index:
    \begin{itemize}
        \item \textbf{Wirtschaft (W)}: Ökonomische Leistungsfähigkeit und Stabilität
        \item \textbf{Gerechtigkeit (G)}: Verteilungsgerechtigkeit und soziale Teilhabe
        \item \textbf{Zukunft (Z)}: Nachhaltigkeit und Investitionen in zukünftiges Wohlergehen
        \item \textbf{Zusammenhalt (S)}: Soziales Kapital und gesellschaftliche Kohäsion
    \end{itemize}
    
    \item[GWI-D-Gewichtung] \hfill \\
    Die relative Bedeutung der vier Komponenten im Gesamtindex: Wirtschaft (0,30), Gerechtigkeit (0,30), Zukunft (0,25), Zusammenhalt (0,15).
    
    \item[GWI-D-Skala] \hfill \\
    Eine normierte Skala von 0 bis 100, wobei 0 den minimalen und 100 den maximalen Gemeinwohlzustand darstellt (theoretische Extremwerte).
    
    \item[GWI-D-Änderung] \hfill \\
    Die Veränderung des Index oder seiner Komponenten über einen bestimmten Zeitraum, gemessen in Indexpunkten.
    
    \item[GWI-D-Korrelation] \hfill \\
    Statistischer Zusammenhang zwischen verschiedenen GWI-D-Komponenten oder zwischen GWI-D und anderen Indikatoren (z.B. BIP).
    
    \item[GWI-D-Paradox] \hfill \\
    Situationen, in denen sich GWI-D und traditionelle Wirtschaftsindikatoren (wie BIP) gegenläufig entwickeln.
    
    \item[GWI-D-Kompass] \hfill \\
    Metaphorische Bezeichnung für die Nutzung des GWI-D als Orientierungsrahmen für politische Entscheidungen.
\end{description}

\section{Ökonomische und soziale Indikatoren}

\begin{description}
    \item[BIP (Bruttoinlandsprodukt)] \hfill \\
    Gesamtwert aller Waren und Dienstleistungen, die in einer Volkswirtschaft innerhalb eines Jahres produziert werden. Traditioneller Hauptindikator für wirtschaftliche Leistung.
    
    \item[Gini-Koeffizient] \hfill \\
    Maß für die Ungleichheit der Einkommens- oder Vermögensverteilung. Wert zwischen 0 (vollkommene Gleichheit) und 1 (vollkommene Ungleichheit).
    
    \item[Armutsquote] \hfill \\
    Anteil der Bevölkerung mit einem Einkommen unterhalb der Armutsrisikoschwelle (meist 60\% des Medianeinkommens).
    
    \item[Gender Pay Gap] \hfill \\
    Prozentualer Unterschied zwischen dem durchschnittlichen Bruttostundenverdienst von Frauen und Männern.
    
    \item[Arbeitsproduktivität] \hfill \\
    Wirtschaftliche Leistung pro Arbeitsstunde oder pro Erwerbstätigen.
    
    \item[Investitionsquote] \hfill \\
    Anteil der Bruttoanlageinvestitionen am Bruttoinlandsprodukt.
    
    \item[Exportquote] \hfill \\
    Anteil der Exporte am Bruttoinlandsprodukt.
    
    \item[Forschungsquote] \hfill \\
    Anteil der Ausgaben für Forschung und Entwicklung am Bruttoinlandsprodukt.
    
    \item[CO2-Intensität] \hfill \\
    Ausstoß von Kohlendioxid pro Einheit des Bruttoinlandsprodukts oder pro Kopf.
\end{description}

\section{Statistische und methodische Begriffe}

\begin{description}
    \item[Korrelation (r)] \hfill \\
    Statistisches Maß für den linearen Zusammenhang zwischen zwei Variablen. Werte zwischen -1 (perfekte negative Korrelation) und +1 (perfekte positive Korrelation).
    
    \item[Signifikanz (p-Wert)] \hfill \\
    Wahrscheinlichkeit, dass ein beobachteter Zusammenhang zufällig auftritt. Übliches Signifikanzniveau: \(p < 0,05\).
    
    \item[Normalisierung] \hfill \\
    Transformation von Daten auf eine einheitliche Skala (im GWI-D: 0–100).
    
    \item[Aggregation] \hfill \\
    Zusammenfassung mehrerer Indikatoren zu einem Gesamtindex.
    
    \item[Zeitreihenanalyse] \hfill \\
    Statistische Untersuchung von Daten über einen Zeitverlauf.
    
    \item[Quantilregression] \hfill \\
    Regressionsmethode, die nicht nur den durchschnittlichen Zusammenhang, sondern Zusammenhänge in verschiedenen Bereichen der Verteilung schätzt.
    
    \item[Imputation] \hfill \\
    Verfahren zur Schätzung fehlender Werte in Datensätzen.
    
    \item[Sensitivitätsanalyse] \hfill \\
    Untersuchung, wie robust Ergebnisse gegenüber methodischen Änderungen (z.B. Gewichtungen) sind.
\end{description}

\section{Historische Perioden und Epochen}

\begin{description}
    \item[Wirtschaftswunder (1960–1973)] \hfill \\
    Phase außergewöhnlich hohen Wirtschaftswachstums und wirtschaftlichen Aufbaus in Westdeutschland nach dem Zweiten Weltkrieg.
    
    \item[Stagflation (1974–1982)] \hfill \\
    Periode wirtschaftlicher Stagnation bei gleichzeitig hoher Inflation, ausgelöst durch Ölpreisschocks und das Ende des Bretton-Woods-Systems.
    
    \item[Konsolidierung (1983–1998)] \hfill \\
    Ära unter Bundeskanzler Helmut Kohl mit wirtschaftlicher Stabilisierung, aber auch sozialen Kürzungen und der deutschen Einheit.
    
    \item[Agenda-Jahre (1999–2008)] \hfill \\
    Periode der Agenda-2010-Reformen unter der rot-grünen Bundesregierung mit tiefgreifenden Veränderungen in Arbeitsmarkt und Sozialstaat.
    
    \item[Krisenbewältigung (2009–2020)] \hfill \\
    Dekade multipler Krisen: Finanzkrise, Eurokrise, Flüchtlingskrise, beginnende Pandemie.
    
    \item[Multikrise (2021–2023)] \hfill \\
    Gleichzeitige Überlagerung von Pandemie, Lieferkettenkrise, Energiekrise, Inflation und geopolitischen Spannungen.
    
    \item[Deutsche Einheit] \hfill \\
    Der Prozess der Wiedervereinigung der Bundesrepublik Deutschland und der Deutschen Demokratischen Republik 1990.
\end{description}

\section{Politische und institutionelle Begriffe}

\begin{description}
    \item[Soziale Marktwirtschaft] \hfill \\
    Wirtschaftsordnung, die Marktwirtschaft mit sozialem Ausgleich verbindet. Gründungsidee der Bundesrepublik Deutschland.
    
    \item[Neoliberalismus] \hfill \\
    Wirtschaftspolitische Ausrichtung, die Marktliberalisierung, Deregulierung und Privatisierung betont und staatliche Interventionen zurückdrängt.
    
    \item[Keynesianismus] \hfill \\
    Wirtschaftstheorie, die aktive staatliche Konjunkturpolitik (insbesondere in Rezessionen) befürwortet.
    
    \item[Schwarze Null] \hfill \\
    Haushaltspolitisches Ziel eines ausgeglichenen Staatshaushalts ohne neue Schulden.
    
    \item[Schuldenbremse] \hfill \\
    Verfassungsrechtliche Regelung, die die Nettokreditaufnahme des Bundes auf 0,35\% des BIP begrenzt.
    
    \item[Sozialpartnerschaft] \hfill \\
    Kooperatives Verhältnis zwischen Gewerkschaften und Arbeitgeberverbänden, charakteristisch für das deutsche Arbeitsmodell.
    
    \item[Mitbestimmung] \hfill \\
    Beteiligung von Arbeitnehmervertretern an betrieblichen und unternehmerischen Entscheidungen.
    
    \item[Bürgergeld] \hfill \\
    Grundsicherungsleistung für erwerbsfähige Personen, eingeführt 2023 als Nachfolger von Hartz IV.
    
    \item[Energiewende] \hfill \\
    Umstieg von fossilen und nuklearen Energieträgern auf erneuerbare Energien.
\end{description}

\section{Abkürzungen und Akronyme}

\begin{description}
    \item[GWI-D] Gemeinwohl-Index Deutschland
    \item[BIP] Bruttoinlandsprodukt
    \item[SOEP] Sozio-oekonomisches Panel
    \item[IAB] Institut für Arbeitsmarkt- und Berufsforschung
    \item[BMBF] Bundesministerium für Bildung und Forschung
    \item[UBA] Umweltbundesamt
    \item[OECD] Organisation für wirtschaftliche Zusammenarbeit und Entwicklung
    \item[EU] Europäische Union
    \item[GDPR] General Data Protection Regulation (Datenschutz-Grundverordnung)
    \item[ALG] Arbeitslosengeld
    \item[ALG II] Arbeitslosengeld II (vormals Hartz IV)
    \item[HDI] Human Development Index (Index der menschlichen Entwicklung)
    \item[SDG] Sustainable Development Goals (Nachhaltigkeitsziele der UN)
    \item[ESS] European Social Survey (Europäische Sozialstudie)
    \item[ALLBUS] Allgemeine Bevölkerungsumfrage der Sozialwissenschaften
    \item[VAR] Vector Autoregression (Vektorautoregression)
    \item[ARIMA] Autoregressive Integrated Moving Average (Autoregressiver integrierter gleitender Durchschnitt)
    \item[CO2] Kohlendioxid
    \item[EU-ETS] European Union Emissions Trading System (EU-Emissionshandelssystem)
    \item[KfW] Kreditanstalt für Wiederaufbau
    \item[BMWK] Bundesministerium für Wirtschaft und Klimaschutz
    \item[BMVI] Bundesministerium für Verkehr und digitale Infrastruktur
    \item[BKA] Bundeskriminalamt
    \item[PKS] Polizeiliche Kriminalstatistik
    \item[GKV] Gesetzliche Krankenversicherung
    \item[PKV] Private Krankenversicherung
    \item[BBSR] Bundesinstitut für Bau-, Stadt- und Raumforschung
    \item[BAMF] Bundesamt für Migration und Flüchtlinge
    \item[AfD] Alternative für Deutschland
    \item[RAF] Rote Armee Fraktion
    \item[ESM] Europäischer Stabilitätsmechanismus
    \item[COVID-19] Coronavirus Disease 2019
    \item[WHO] Weltgesundheitsorganisation
    \item[IMF] Internationaler Währungsfonds
    \item[ECB] Europäische Zentralbank
\end{description}

\section{Maßeinheiten und Skalen}

\begin{description}
    \item[Indexpunkte] Einheiten, in denen Veränderungen im GWI-D gemessen werden.
    \item[Prozentpunkte (PP)] Differenz zwischen zwei Prozentwerten.
    \item[Mrd. Euro] Milliarden Euro (10.000.000.000 Euro).
    \item[t/Kopf] Tonnen pro Kopf.
    \item[\% BIP] Prozent des Bruttoinlandsprodukts.
    \item[€/Monat] Euro pro Monat.
    \item[Stunden/Woche] Stunden pro Woche.
\end{description}

\section*{Hinweise zur Nutzung}

\begin{itemize}
    \item Dieses Glossar erhebt keinen Anspruch auf Vollständigkeit, sondern konzentriert sich auf die in der GWI-D-Analyse zentralen Begriffe.
    \item Bei weiteren Fragen zur Terminologie oder Methodik wird auf die entsprechenden Kapitel der Hauptanalyse verwiesen.
    \item Das Glossar wird regelmäßig aktualisiert und erweitert, insbesondere bei Weiterentwicklung des GWI-D-Konzepts.
    \item Für spezifische statistische oder methodische Fragen wird auf die Fachliteratur und die detaillierte Methodikdokumentation verwiesen.
\end{itemize}

\endinput